\subsection{Zero Knowledge}
\label{defn:zk}

We next turn to defining zero-knowledge for a DKG.
Zero knowledge requires that a probabilistic polynomial-time adversary $\adv$ allowed to control up to $t-1$ participants does not gain any additional advantage in its effort to learn $\SK$ than it would against the target key generation algorithm.
This notion is formalized by the requirement of \emph{simulatability}:
that the DKG can be simulated to an adversary with respect to some challenge $\dkgchal \in \publicspace$,
where $\dkgchal$ is a valid public key for the target key generation algorithm.
We formalize this notion in Figure~\ref{fg:zk}.
The oracle $\oraclekg$ is defined in Figure~\ref{fg:keygen-o}.

Note that we implicitly assume weak robustness,
as the adversary cannot win if it causes all participants to non-fatally abort.

\begin{figure}[!ht]
  \begin{pchstack}[boxed,center,space=1em]
  \begin{pcvstack}
    \procedure[]{$\ZKExp_{\dkg,\adv,\simkeygen_i,\simfin}(\lambda)$}{
    b \rgets \zo \\
    \randcoin \rgets \zo \text{\algcom Random coin for losing conditions} \\
    \roundcounter \gets 0 \text{\algcom Ensure rounds are queried in order}\\
    \invecp_{11} \gets \emptyset, \ldots, \invecp_{1n} \gets \emptyset \text{\algcom Simulate peer-to-peer channels} \\
    (n, t, \corrupt, \advstate) \gets \adv(1^\lambda);\
    \honest \gets \set{n} \setminus \corrupt \\
    \pcif n < t \vee \corrupt \not\subset \set{n} \vee \lvert \corrupt \rvert \geq t \\
    \pcind \pcreturn \randcoin \\
    (\cdot, \dkgchal) \rgets \targetkeygen.\keygen(1^\lambda)  \\
%
    \pcfor i \in \honest \pcdo \\
    \pcind \pcif b=0 \pcthen (\party_i, \outvecp_i, \bmsg_i) \gets \dkgkeygen_0(\lambda, n, t, i)  \\
    \pcind \pcelse (\party_i, \outvecp_i, \bmsg_i) \gets \simkeygen_0[\Hash](\lambda, n, t, i, \dkgchal)  \\
    \pcind \invecp_{1k}[i] \gets \outvecp_{0i}[k],\ \forall\ k \in \set{n} \\
    \text{\algcom Set party $i$'s peer-to-peer messages for all other parties} \\ % Make sure all parties receive peer-to-peer messages
%
    b' \gets \adv^{\oraclekg,\oraclefinalize}(\advstate, \{\invecp_{1j} \}_{j \in \corrupt}, \{ \bmsg_i \}_{i \in \honest} ) \\
    \pcif \roundcounter \neq \numrounds \text{\algcom Prevents trivial win by forcing coordination} \\
    \pcind \pcreturn \randcoin \\
    \pcif \status_j = \abort \text{ for all } j \in \honest \\
    \pcind \pcreturn \randcoin \\
    \text{\algcom $\adv$ cannot win if all honest parties non-fatally abort} \\
    %
    \pcif b=1 \\
    \pcind \pcreturn 1\ \pcif \exists\ i \in \honest :  \PK_i \neq \dkgchal \\
    \text{\algcom $\adv$ wins when $b=1$ if the game is not simulated with respect to $\dkgchal$} \\
    %
    \pcreturn 1\ \pcif b \iseq b' \\
    \pcreturn 0
  }
  \end{pcvstack}
\end{pchstack}

\caption{Zero-knowledge experiment for a DKG.
  }
\label{fg:zk}
\end{figure}


Our definition of zero-knowledge is distinct from the simulation-based notion of secrecy
introduced by Gennaro et al.~\cite{GennaroJKR07} in that we formally define the
notion with respect to a target key generation algorithm.
whereas the definition by Gennaro et al.\ assumes random group elements
as the public key to simulate.
In some settings (where public and secret keys have a unique correspondence),
then these definitions are the same.
If not, then such simulation-based notions fail to consider
indistinguishability, as discussed below.

We now define this notion of zero-knowledge more formally.
The advantage of an adversary $\adv$ against a DKG $\dkg$
with respect to the simulation algorithms $(\simkeygen_i, i \in \{0,\ldots,\numrounds-1 \}, \simfin)$
in the zero-knowledge experiment as defined in Figure~\ref{fg:zk}
is
\[
  \advantzk{\dkg,\adv,\simkeygen_i,\simfin}(\lambda) = \bigabs{\Pr[\ZKExp_{\dkg,\adv,\simkeygen_i,\simfin}(\lambda) = 1 ] - 1/2   }
\]

\begin{definition}
  A DKG $\dkg$ is \defn{zero-knowledge} if there exists efficient algorithms $(\simkeygen_i, i \in \{0,\ldots,\numrounds-1 \}, \simfin)$ such that
  for all probabilistic polynomial-time adversaries $\adv$,
  $\advantzk{\dkg,\adv,\simkeygen_i,\simfin}(\lambda)$ is negligible.
\end{definition}

The zero-knowledge attack game in Figure~\ref{fg:zk} requires $\adv$ to successfully guess if the environment that it is in is real or simulated with respect to a challenge $\dkgchal$,
where $\dkgchal$ is a public key generated by $\targetkeygen.\keygen$.
The game accepts as input the security parameter $\lambda$,
and internally to the experiment, samples a bit $b \in \zo$ at random.
In the simulated setting (when $b=1$),
the environment simulates the honest participants with respect to $\dkgchal$ by
employing the simulation algorithms $(\simkeygen_i, i \in \{0,\ldots,\numrounds-1 \}, \simfin)$.
The adversary wins by default if the environment cannot simulate the protocol;
i.e, the resulting group public key $\pk$ does not equal $\dkgchal$.
Otherwise, the adversary wins if it successfully guesses whether it is playing against the real or simulated protocol.


